\documentclass[10pt,a4paper]{article}
\usepackage[margin=15mm,top=12mm,bottom=12mm]{geometry}
\usepackage{fontspec}
\usepackage[french]{babel}
\usepackage{amsmath,amssymb}
\usepackage{xcolor}
\usepackage[most]{tcolorbox}
\usepackage{booktabs}
\usepackage{enumitem}
\usepackage{graphicx}
\usepackage{array}
\usepackage{titlesec}
\usepackage{listings}
\usepackage{hyperref}

\definecolor{bleu}{HTML}{1F6FEB}
\definecolor{gris}{HTML}{5B6470}
\definecolor{vert}{HTML}{1E8449}
\definecolor{orange}{HTML}{B45309}
\definecolor{fond}{HTML}{F4F6F8}
\definecolor{codebg}{HTML}{F6F8FA}

\setlength{\parindent}{0pt}
\titlespacing*{\section}{0pt}{10pt}{4pt}
\titleformat{\section}{\normalfont\bfseries\color{bleu}\large}{}{0pt}{}
\titleformat{\subsection}{\normalfont\bfseries\normalsize}{}{0pt}{}
\titlespacing*{\subsection}{0pt}{8pt}{3pt}
\setlist[itemize]{nosep,leftmargin=1.1em,labelsep=.45em}
\setlist[enumerate]{nosep,leftmargin=1.35em,labelsep=.5em}

\lstdefinestyle{py}{language=Python,backgroundcolor=\color{codebg},
  basicstyle=\ttfamily\fontsize{7.6}{9}\selectfont,
  keywordstyle=\color{bleu}\bfseries,commentstyle=\color{vert}\itshape,
  stringstyle=\color{orange},frame=none,breaklines=true,
  showstringspaces=false,columns=fullflexible,keepspaces=true,
  aboveskip=2pt,belowskip=2pt,
  literate={é}{{\'e}}1 {è}{{\`e}}1 {à}{{\`a}}1 {ç}{{\c c}}1 {û}{{\^u}}1}
\lstset{style=py}

\newtcolorbox{corrige}{enhanced,colback=vert!4,colframe=vert!60!black,boxrule=.5pt,
  arc=2pt,left=5pt,right=5pt,top=3pt,bottom=3pt,breakable,
  fonttitle=\bfseries\footnotesize\color{white},colbacktitle=vert!70!black,
  title={Éléments de réponse --- à ne lire qu'après avoir cherché}}
\newtcolorbox{enonce}{enhanced,colback=fond,colframe=fond,arc=2pt,
  left=5pt,right=5pt,top=3pt,bottom=3pt,breakable}

\begin{document}

\begin{tcolorbox}[enhanced,colback=fond,colframe=fond,arc=3pt,
  left=9pt,right=9pt,top=7pt,bottom=7pt]
\begin{minipage}[c]{0.76\textwidth}
  {\color{gris}\footnotesize LICENCE 3 SDIA $\cdot$ ENSPD --- UNIVERSITÉ DE DOUALA}\\[2pt]
  {\Large\bfseries Travaux dirigés --- recueil complet}\\[3pt]
  {\small Introduction au Machine Learning $\cdot$ 8 chapitres $\cdot$ 76 exercices}\\[4pt]
  {\color{gris}\footnotesize Dr.\ Noulapeu Ngaffo Armielle}
\end{minipage}\hfill
\begin{minipage}[c]{0.20\textwidth}
  \raggedleft\includegraphics[width=17mm]{logo_enspd.png}
\end{minipage}
\end{tcolorbox}

\vspace{2mm}
{\small
\textbf{Mode d'emploi.} Chaque TD suit la progression du cours : calculs à la main,
exercice de code, diagnostic, question de synthèse. \textbf{Aucun corrigé n'est
inclus} : ce recueil est destiné à la distribution en séance. Sauf mention contraire,
les calculs se font sans calculatrice.}

\vspace{1.5mm}
\begin{tcolorbox}[enhanced,colback=fond,colframe=fond,arc=2pt,left=6pt,right=6pt,
  top=4pt,bottom=4pt]
{\small
\textbf{Organisation conseillée.} Chaque TD correspond à un chapitre et se traite
en \textbf{2 h} : 30 min de recherche individuelle, 50 min de mise en commun,
40 min de correction au tableau. Les neuf exercices sont progressifs : les quatre
premiers sont faisables en séance, les suivants demandent un travail personnel.\\[2pt]
\textbf{Volume total} : 76 exercices, soit environ 16 h de travail étudiant, et un
TD de synthèse qui prépare le projet intégrateur. Les \textbf{TD ne
remplacent pas les TP} : les TP font exécuter du code, les TD font raisonner et
calculer à la main.}
\end{tcolorbox}

\vspace{1mm}
\tableofcontents
\vspace{2mm}
\newpage

\section*{TD 1 --- Démarche et généralisation}
\addcontentsline{toc}{section}{TD 1 --- Démarche et généralisation}

\subsection*{Exercice 1 --- Vocabulaire \hfill \normalfont\footnotesize (5 min)}
Pour chaque situation, donner $X$ et $y$, puis le type de tâche.
\begin{enumerate}
  \item Prédire la consommation électrique d'un bâtiment à partir de l'heure,
        du jour et de la température.
  \item Regrouper des clients selon leurs achats, sans catégorie connue.
  \item Décider si un dossier de crédit est risqué, à partir de revenus et
        d'historiques de remboursement.
\end{enumerate}

\subsection*{Exercice 2 --- Risque empirique \hfill \normalfont\footnotesize (10 min)}
On dispose de quatre observations :

\begin{center}
\small
\begin{tabular}{@{}lrrrr@{}}
\toprule
$i$ & 1 & 2 & 3 & 4 \\
\midrule
$y_i$ & 100 & 150 & 80 & 200 \\
$\hat y_i$ & 90 & 120 & 100 & 180 \\
\bottomrule
\end{tabular}
\end{center}

\begin{enumerate}
  \item Calculer la perte quadratique moyenne $\hat R_4$ et la RMSE.
  \item Calculer la perte absolue moyenne.
  \item Une erreur de 30 pèse combien de fois plus qu'une erreur de 10 dans
        chaque perte ?
\end{enumerate}

\subsection*{Exercice 3 --- Diagnostiquer \hfill \normalfont\footnotesize (8 min)}
Trois modèles, trois bilans. Pour chacun : sous-apprentissage, bon compromis ou
surapprentissage ? Justifier en une phrase.

\begin{center}
\small
\begin{tabular}{@{}lccc@{}}
\toprule
& \textbf{Erreur entraînement} & \textbf{Erreur test} & \textbf{Diagnostic} \\
\midrule
A & 0,32 & 0,34 & \\
B & 0,02 & 0,41 & \\
C & 0,08 & 0,09 & \\
\bottomrule
\end{tabular}
\end{center}

\subsection*{Exercice 4 --- Choisir le découpage \hfill \normalfont\footnotesize (8 min)}
Quel découpage (aléatoire, temporel, groupé, stratifié) imposez-vous ?
\begin{enumerate}
  \item Prédire les ventes du mois suivant à partir de 3 ans d'historique.
  \item Détecter une maladie rare : 4 \% de cas positifs.
  \item Classer des élèves dont certains ont plusieurs relevés annuels.
\end{enumerate}

\subsection*{Exercice 5 --- Fuite ou pas ? \hfill \normalfont\footnotesize (8 min)}
Un groupe prédit la réussite d'étudiants en octobre. Son notebook contient :
\begin{itemize}
  \item « moyenne du premier semestre » ;
  \item « nombre d'absences » ;
  \item « a été convoqué en tutorat en juin » ;
  \item une standardisation calculée sur l'ensemble des données avant le split.
\end{itemize}
Identifier les deux problèmes et proposer une correction pour chacun.


\subsection*{Exercice 6 --- Exactitude sur huit dossiers \hfill \normalfont\footnotesize (6 min)}
Un modèle classe 8 dossiers : 6 sont correctement prédits.
\begin{enumerate}
  \item Calculer l'exactitude.
  \item Ce chiffre suffit-il si 7 dossiers sur 8 sont négatifs ? Justifier.
  \item Quelle information manque pour conclure ?
\end{enumerate}

\subsection*{Exercice 7 --- Choisir la perte \hfill \normalfont\footnotesize (6 min)}
On prédit le \textbf{nombre d'incidents} mensuels d'un serveur (comptage).
\begin{enumerate}
  \item Pourquoi la perte quadratique est-elle peu adaptée ?
  \item Quelle perte privilégier pour un comptage ?
  \item Que se passe-t-il si l'on prédit une valeur négative ?
\end{enumerate}

\subsection*{Exercice 8 --- Le score trop beau \hfill \normalfont\footnotesize (5 min)}
Un groupe annonce 100 \% d'exactitude sur un problème où les données contiennent
du bruit. Rédiger les trois vérifications à demander avant de le croire.

\subsection*{Exercice 9 --- 60/20/20 ou validation croisée ? \hfill \normalfont\footnotesize (6 min)}
Donner un avantage et un inconvénient de chaque approche, puis indiquer laquelle
choisir avec 120 observations.

\newpage

\section*{TD 2 --- Préparation des données}
\addcontentsline{toc}{section}{TD 2 --- Préparation des données}

\subsection*{Exercice 1 --- Que faire des manquants ? \hfill \normalfont\footnotesize (8 min)}
Une table de 1\,000 étudiants présente 41 \% de valeurs manquantes sur le
revenu des parents, et 3 \% sur la moyenne générale.
\begin{enumerate}
  \item Proposer une décision pour chaque colonne, en justifiant.
  \item Pourquoi supprimer les lignes à 41 \% de manquants est-il risqué ?
  \item Quelle information ajouter pour améliorer l'imputation ?
\end{enumerate}

\subsection*{Exercice 2 --- Encodages \hfill \normalfont\footnotesize (8 min)}
Trois colonnes : \texttt{filiere} (informatique, télécom, réseaux),
\texttt{niveau} (faible, moyen, élevé), \texttt{code\_postal} (400 modalités).
Indiquer l'encodage retenu pour chacune et la raison.

\subsection*{Exercice 3 --- Effet de l'échelle \hfill \normalfont\footnotesize (10 min)}
Deux étudiants sont décrits par (ancienneté en années, revenu en FCFA) :
$A=(0,0)$ et $B=(3,\,4\,000)$.
\begin{enumerate}
  \item Calculer la distance euclidienne brute.
  \item Les écarts-types valent 1,5 an et 2\,000 FCFA. Recalculer la distance
        après standardisation.
  \item Que conclure sur l'usage de kNN sans normalisation ?
\end{enumerate}

\subsection*{Exercice 4 --- Détecter une valeur aberrante \hfill \normalfont\footnotesize (8 min)}
Série : $10,\,12,\,14,\,16,\,18,\,20,\,22,\,60$. On utilise la méthode de Tukey
(médiane de chaque moitié).
\begin{enumerate}
  \item Calculer $Q_1$, $Q_3$, l'IQR et les deux bornes.
  \item La valeur 60 est-elle aberrante ? Et la valeur 10 ?
  \item Faut-il la supprimer ? Que faut-il vérifier avant ?
\end{enumerate}

\subsection*{Exercice 5 --- Trouver la fuite dans le code \hfill \normalfont\footnotesize (10 min)}
\begin{lstlisting}
X_scaled = StandardScaler().fit_transform(X)
X_train, X_test, y_train, y_test = train_test_split(X_scaled, y, test_size=0.2)
imputer = SimpleImputer(strategy="mean").fit(X)
X_train = imputer.transform(X_train)
\end{lstlisting}
Deux lignes posent problème. Les identifier et réécrire le pipeline correctement.


\subsection*{Exercice 6 --- Encodage cyclique \hfill \normalfont\footnotesize (8 min)}
L'heure est encodée par $h_{\sin}=\sin(2\pi h/24)$ et $h_{\cos}=\cos(2\pi h/24)$.
\begin{enumerate}
  \item Calculer le couple $(\sin,\cos)$ pour 0 h, 6 h, 12 h et 18 h.
  \item Quelle est la distance euclidienne entre 23 h et 0 h ? Entre 0 h et 12 h ?
  \item Pourquoi un encodage entier de l'heure est-il trompeur ?
\end{enumerate}

\subsection*{Exercice 7 --- Imputation et dispersion \hfill \normalfont\footnotesize (7 min)}
On impute 25 \% de valeurs manquantes d'un revenu très asymétrique.
\begin{enumerate}
  \item L'imputation par la moyenne conserve-t-elle la moyenne ? La variance ?
  \item Que fait l'imputation par la médiane à la moyenne globale ?
  \item Quelle information perd-on dans les deux cas ?
\end{enumerate}

\subsection*{Exercice 8 --- Discrétiser ou non ? \hfill \normalfont\footnotesize (6 min)}
On découpe l'âge en quatre tranches.
\begin{enumerate}
  \item Citer deux avantages de ce choix.
  \item Citer deux inconvénients.
  \item Pour quel type de modèle la discrétisation est-elle la moins utile ?
\end{enumerate}

\subsection*{Exercice 9 --- La jointure qui double les lignes \hfill \normalfont\footnotesize (8 min)}
Une jointure fait passer de 400 étudiants à 752 lignes.
\begin{enumerate}
  \item Expliquer la cause probable.
  \item Quel effet sur le taux de réussite moyen ?
  \item Proposer deux corrections possibles.
\end{enumerate}

\newpage

\section*{TD 3 --- Régression et régularisation}
\addcontentsline{toc}{section}{TD 3 --- Régression et régularisation}

\subsection*{Exercice 1 --- Ajuster une droite à la main \hfill \normalfont\footnotesize (12 min)}
Quatre points : $(1,3)$, $(2,5)$, $(3,6)$, $(4,8)$.
\begin{enumerate}
  \item Calculer $\bar x$ et $\bar y$.
  \item Calculer la pente $\hat\beta_1=\dfrac{\sum (x_i-\bar x)(y_i-\bar y)}{\sum (x_i-\bar x)^2}$.
  \item En déduire l'ordonnée à l'origine, puis la prédiction en $x=5$.
  \item Calculer le $R^2$.
\end{enumerate}

\subsection*{Exercice 2 --- Interpréter un coefficient \hfill \normalfont\footnotesize (6 min)}
Un modèle estime $\hat y = 15{,}2 + 0{,}42\,x$ où $x$ est la surface en m$^2$
et $y$ le loyer en milliers de FCFA.
\begin{enumerate}
  \item Que signifie le coefficient $0{,}42$ ?
  \item Que signifie l'ordonnée $15{,}2$ ? A-t-elle un sens ici ?
  \item Prédire le loyer pour 60 m$^2$.
\end{enumerate}

\subsection*{Exercice 3 --- Ridge ou Lasso ? \hfill \normalfont\footnotesize (8 min)}
On dispose de 40 variables, dont beaucoup sont redondantes.
\begin{enumerate}
  \item Quelle pénalité privilégier si l'on veut \textbf{sélectionner} des
        variables ? Pourquoi ?
  \item Laquelle privilégier si l'on veut seulement \textbf{stabiliser} les
        coefficients sans en annuler ?
  \item Comment choisit-on $\lambda$ ?
\end{enumerate}

\subsection*{Exercice 4 --- RMSE contre MAE \hfill \normalfont\footnotesize (8 min)}
Deux modèles, sur le même jeu de test :
\begin{center}
\small
\begin{tabular}{@{}lcc@{}}
\toprule
& \textbf{RMSE} & \textbf{MAE} \\
\midrule
Modèle A & 12,0 & 11,4 \\
Modèle B & 15,2 & 8,1 \\
\bottomrule
\end{tabular}
\end{center}
\begin{enumerate}
  \item Lequel a les plus grosses erreurs ponctuelles ?
  \item Lequel choisir si une grosse erreur coûte très cher ?
  \item Lequel choisir si l'on veut l'erreur typique la plus faible ?
\end{enumerate}

\subsection*{Exercice 5 --- Diagnostic de résidus \hfill \normalfont\footnotesize (8 min)}
On observe un nuage de résidus en forme de U en fonction de la prédiction.
\begin{enumerate}
  \item Quel défaut du modèle cela révèle-t-il ?
  \item Citer deux remèdes possibles.
  \item Que signifierait au contraire un entonnoir (variance croissante) ?
\end{enumerate}


\subsection*{Exercice 6 --- Interpréter un $R^2$ \hfill \normalfont\footnotesize (6 min)}
Deux modèles obtiennent $R^2=0{,}98$ et $R^2=0{,}45$ sur le même problème.
\begin{enumerate}
  \item Que signifie chaque valeur en termes de variance expliquée ?
  \item Le $R^2$ élevé garantit-il de bonnes prévisions ? Pourquoi ?
  \item Que ne dit pas le $R^2$ ?
\end{enumerate}

\subsection*{Exercice 7 --- Multicolinéarité \hfill \normalfont\footnotesize (7 min)}
Deux variables sont corrélées à $0{,}95$ et leurs coefficients changent de signe
d'un échantillon à l'autre.
\begin{enumerate}
  \item Nommer le phénomène.
  \item Pourquoi les coefficients deviennent-ils instables ?
  \item Citer deux remèdes.
\end{enumerate}

\subsection*{Exercice 8 --- Une descente de gradient à la main \hfill \normalfont\footnotesize (10 min)}
On minimise $J(w)=(w-3)^2+1$, avec $w_0=0$ et $\eta=0{,}1$.
\begin{enumerate}
  \item Calculer le gradient en $w_0$.
  \item En déduire $w_1$, puis $w_2$ et $w_3$.
  \item Vers quelle valeur la suite converge-t-elle ?
\end{enumerate}

\subsection*{Exercice 9 --- MAE ou MSE ? \hfill \normalfont\footnotesize (6 min)}
Les données de prix contiennent quelques valeurs extrêmes.
\begin{enumerate}
  \item Quelle perte est la plus robuste ? Pourquoi ?
  \item Quelle perte pénalise davantage une grosse erreur ponctuelle ?
  \item Laquelle choisir si le coût est proportionnel à l'erreur ?
\end{enumerate}

\newpage

\section*{TD 4 --- Classification et évaluation}
\addcontentsline{toc}{section}{TD 4 --- Classification et évaluation}

\subsection*{Exercice 1 --- Lire une matrice de confusion \hfill \normalfont\footnotesize (12 min)}
Sur 200 dossiers :

\begin{center}
\small
\begin{tabular}{@{}lrr@{}}
\toprule
& \textbf{prédit sain} & \textbf{prédit malade} \\
\midrule
\textbf{réel sain} & 130 & 10 \\
\textbf{réel malade} & 20 & 40 \\
\bottomrule
\end{tabular}
\end{center}

\begin{enumerate}
  \item Calculer l'exactitude, la précision, le rappel et $F_1$.
  \item Que signifie concrètement « précision = 0,80 » ?
  \item Le modèle rate 20 malades sur 60. Est-ce acceptable ? De quoi dépend la
        réponse ?
\end{enumerate}

\subsection*{Exercice 2 --- Déplacer le seuil \hfill \normalfont\footnotesize (8 min)}
Le modèle produit une probabilité. On abaisse le seuil de 0,5 à 0,3.
\begin{enumerate}
  \item Le rappel augmente-t-il ou diminue-t-il ?
  \item Et la précision ?
  \item Dans quel contexte ce choix est-il justifié ?
\end{enumerate}

\subsection*{Exercice 3 --- ROC et AUC \hfill \normalfont\footnotesize (8 min)}
\begin{enumerate}
  \item Que représente un point de la courbe ROC ?
  \item Que vaut l'AUC d'un classifieur aléatoire ? D'un classifieur parfait ?
  \item Pourquoi l'AUC peut-elle être trompeuse sur des classes très
        déséquilibrées ?
\end{enumerate}

\subsection*{Exercice 4 --- Le piège de l'exactitude \hfill \normalfont\footnotesize (8 min)}
Sur 2\,000 dossiers, 5 \% sont positifs. Un modèle prédit toujours « négatif ».
\begin{enumerate}
  \item Calculer son exactitude et son rappel.
  \item Pourquoi ce modèle est-il inutile malgré 95 \% ?
  \item Quelle métrique publier à la place ?
\end{enumerate}

\subsection*{Exercice 5 --- Choisir la métrique \hfill \normalfont\footnotesize (6 min)}
Indiquer la métrique à privilégier et pourquoi :
\begin{enumerate}
  \item dépistage d'une maladie rare ;
  \item filtre anti-spam ;
  \item prévision de stock (erreur en unités).
\end{enumerate}


\subsection*{Exercice 6 --- Précision-rappel ou ROC ? \hfill \normalfont\footnotesize (7 min)}
\begin{enumerate}
  \item Que montre une courbe précision-rappel que la courbe ROC ne montre pas ?
  \item Dans quel cas la ROC est-elle trop optimiste ?
  \item Laquelle publier pour un dépistage avec 2 \% de positifs ?
\end{enumerate}

\subsection*{Exercice 7 --- Calibration \hfill \normalfont\footnotesize (7 min)}
Un modèle annonce une probabilité moyenne de 0,90 sur un lot de dossiers, dont
seulement 60 \% sont réellement positifs.
\begin{enumerate}
  \item Le modèle est-il bien classé ? Est-il bien calibré ?
  \item Que signifie concrètement ce biais pour l'utilisateur ?
  \item Comment corriger la calibration ?
\end{enumerate}

\subsection*{Exercice 8 --- Trois classes \hfill \normalfont\footnotesize (10 min)}
Matrice de confusion (150 dossiers, 50 par classe) :

\begin{center}
\small
\begin{tabular}{@{}lrrr@{}}
\toprule
 & \textbf{prédit A} & \textbf{prédit B} & \textbf{prédit C} \\
\midrule
\textbf{réel A} & 40 & 5 & 5 \\
\textbf{réel B} & 10 & 30 & 10 \\
\textbf{réel C} & 5 & 5 & 40 \\
\bottomrule
\end{tabular}
\end{center}

\begin{enumerate}
  \item Calculer l'exactitude globale.
  \item Calculer le rappel de chaque classe.
  \item Quelle classe est la plus difficile, et que faire ensuite ?
\end{enumerate}

\subsection*{Exercice 9 --- Le seuil a un coût \hfill \normalfont\footnotesize (8 min)}
Une fausse alerte coûte 1 unité ; un cas manqué coûte 10 unités.
\begin{enumerate}
  \item À partir de quelle probabilité faut-il alerter ?
  \item Que devient ce seuil si un cas manqué coûte 100 unités ?
  \item Pourquoi le seuil 0,5 n'a-t-il rien de naturel ?
\end{enumerate}

\newpage

\section*{TD 5 --- Arbres et ensembles}
\addcontentsline{toc}{section}{TD 5 --- Arbres et ensembles}

\subsection*{Exercice 1 --- Impureté de Gini \hfill \normalfont\footnotesize (12 min)}
Un nœud contient 10 exemples : 6 positifs, 4 négatifs.
\begin{enumerate}
  \item Calculer son indice de Gini.
  \item On le coupe en deux enfants : $A$ (5 positifs, 1 négatif) et $B$
        (1 positif, 3 négatifs). Calculer le Gini de chacun.
  \item Calculer le Gini pondéré des enfants, puis le gain.
  \item La coupe est-elle utile ?
\end{enumerate}

\subsection*{Exercice 2 --- Arbre ou forêt ? \hfill \normalfont\footnotesize (8 min)}
Un arbre unique atteint 0,96 en entraînement et 0,71 en test.
\begin{enumerate}
  \item Quel est le problème ?
  \item Comment une forêt aléatoire le corrige-t-elle ? Expliquer en termes de
        variance.
  \item Pourquoi une forêt ne descend-elle jamais à 0 d'erreur d'entraînement ?
\end{enumerate}

\subsection*{Exercice 3 --- Importance et redondance \hfill \normalfont\footnotesize (8 min)}
Deux colonnes sont corrélées à $r=0{,}95$, et le modèle leur attribue une
importance de 0,20 à chacune, contre 0,45 pour une troisième variable.
\begin{enumerate}
  \item L'importance totale des deux colonnes redondantes est-elle faible ?
  \item Peut-on conclure que chacune compte peu ?
  \item Que faire pour y voir clair ?
\end{enumerate}

\subsection*{Exercice 4 --- Bagging ou boosting ? \hfill \normalfont\footnotesize (8 min)}
Associer chaque affirmation au bon ensemble.
\begin{enumerate}
  \item Les arbres sont entraînés en parallèle, puis moyennés.
  \item Chaque arbre corrige les erreurs du précédent.
  \item Réduit surtout la variance.
  \item Réduit surtout le biais, mais peut surapprendre si l'on ajoute trop
        d'itérations.
\end{enumerate}

\subsection*{Exercice 5 --- Interpréter une décision \hfill \normalfont\footnotesize (6 min)}
Un modèle refuse un crédit. On vous demande d'expliquer la décision à un
client.
\begin{enumerate}
  \item Pourquoi l'importance globale des variables ne suffit-elle pas ?
  \item Que apporte une explication locale (SHAP) ?
  \item Quelle précaution prendre en présentant ces valeurs ?
\end{enumerate}


\subsection*{Exercice 6 --- Entropie contre Gini \hfill \normalfont\footnotesize (8 min)}
Un nœud contient 6 positifs et 4 négatifs.
\begin{enumerate}
  \item Calculer son indice de Gini (déjà vu au TD précédent).
  \item Calculer son entropie $-\sum p_k\log_2 p_k$.
  \item Les deux critères classent-ils les mêmes découpages de la même façon ?
        Que change le choix ?
\end{enumerate}

\subsection*{Exercice 7 --- Profondeur et généralisation \hfill \normalfont\footnotesize (7 min)}
Un arbre donne les résultats suivants :

\begin{center}
\small
\begin{tabular}{@{}lccc@{}}
\toprule
\textbf{Profondeur} & 2 & 5 & 12 \\
\midrule
Erreur entraînement & 0,22 & 0,08 & 0,01 \\
Erreur validation & 0,24 & 0,11 & 0,19 \\
\bottomrule
\end{tabular}
\end{center}

\begin{enumerate}
  \item Quelle profondeur choisir ?
  \item Que se passe-t-il au-delà ?
  \item Comment obtient-on cette courbe sans toucher au test ?
\end{enumerate}

\subsection*{Exercice 8 --- Pourquoi moyenner aide \hfill \normalfont\footnotesize (9 min)}
Cinq arbres ont une variance individuelle $\sigma^2=100$ et une corrélation
moyenne $\rho=0{,}2$.
\begin{enumerate}
  \item Rappeler la variance d'un arbre seul.
  \item Calculer la variance de leur moyenne :
        $\rho\sigma^2+(1-\rho)\sigma^2/5$.
  \item Que se passerait-il avec $\rho=0{,}9$ ? Conclure sur l'intérêt de
        décorréler les arbres.
\end{enumerate}

\subsection*{Exercice 9 --- Régler un boosting \hfill \normalfont\footnotesize (7 min)}
On augmente le nombre d'itérations de 100 à 2\,000.
\begin{enumerate}
  \item Que fait l'erreur d'entraînement ?
  \item Et l'erreur de validation ?
  \item Quel rôle joue le pas d'apprentissage dans ce compromis ?
\end{enumerate}

\newpage

\section*{TD 6 --- SVM et noyaux}
\addcontentsline{toc}{section}{TD 6 --- SVM et noyaux}

\subsection*{Exercice 1 --- Marge et vecteurs de support \hfill \normalfont\footnotesize (8 min)}
Deux classes sont séparables par une droite. Trois points sont sur la marge :
$A$, $B$ et $C$.
\begin{enumerate}
  \item Qu'est-ce qu'un vecteur de support ?
  \item Que devient la frontière si l'on supprime un point très éloigné de la
        marge ?
  \item Et si l'on supprime $A$ ?
\end{enumerate}

\subsection*{Exercice 2 --- Effet de $C$ \hfill \normalfont\footnotesize (8 min)}
On fait varier $C$ de $0{,}01$ à $100$.
\begin{enumerate}
  \item Que se passe-t-il pour $C$ très petit ?
  \item Et pour $C$ très grand ?
  \item Comment choisir $C$ ?
\end{enumerate}

\subsection*{Exercice 3 --- Noyau RBF \hfill \normalfont\footnotesize (10 min)}
Deux points : $x_1=(1,2)$ et $x_2=(3,4)$. Noyau
$K(x_1,x_2)=\exp(-\gamma\lVert x_1-x_2\rVert^2)$.
\begin{enumerate}
  \item Calculer $\lVert x_1-x_2\rVert^2$.
  \item Calculer $K$ pour $\gamma=0{,}1$, puis $\gamma=0{,}5$, puis $\gamma=2$.
  \item Que se passe-t-il quand $\gamma$ augmente ?
\end{enumerate}

\subsection*{Exercice 4 --- SVM ou arbre ? \hfill \normalfont\footnotesize (8 min)}
Indiquer le modèle le plus adapté et pourquoi :
\begin{enumerate}
  \item 80 observations, 500 variables, frontière nette.
  \item 200\,000 lignes, variables catégorielles et manquantes.
  \item Besoin d'expliquer chaque décision à un client.
\end{enumerate}

\subsection*{Exercice 5 --- Que fait le noyau ? \hfill \normalfont\footnotesize (6 min)}
\begin{enumerate}
  \item Pourquoi ne calcule-t-on pas explicitement la transformation
        $\varphi(x)$ dans un espace de grande dimension ?
  \item Que permet l'astuce du noyau ?
  \item Citer deux noyaux usuels.
\end{enumerate}


\subsection*{Exercice 6 --- Calculer une marge \hfill \normalfont\footnotesize (8 min)}
Le séparateur est $w=(1,1)$, $b=0$.
\begin{enumerate}
  \item Calculer $\lVert w\rVert$.
  \item Calculer la marge $2/\lVert w\rVert$.
  \item Calculer la distance du point $(1,1)$ à la droite. Que remarque-t-on ?
\end{enumerate}

\subsection*{Exercice 7 --- Points mal classés et marge souple \hfill \normalfont\footnotesize (7 min)}
Deux points sont du mauvais côté de la frontière.
\begin{enumerate}
  \item Que se passe-t-il si l'on exige une séparation parfaite ?
  \item Que permet la variable ressort (\emph{slack}) ?
  \item Quel paramètre règle le compromis marge / erreurs ?
\end{enumerate}

\subsection*{Exercice 8 --- Noyau polynomial \hfill \normalfont\footnotesize (6 min)}
Noyau $K(x,x')=(1+x\cdot x')^2$, avec $x=(1,2)$ et $x'=(2,1)$.
\begin{enumerate}
  \item Calculer $x\cdot x'$.
  \item En déduire $K(x,x')$.
  \item Pourquoi ne calcule-t-on pas explicitement $\varphi(x)$ ?
\end{enumerate}

\subsection*{Exercice 9 --- Complexité \hfill \normalfont\footnotesize (6 min)}
Un SVM exact a un coût de l'ordre de $O(n^2)$ à $O(n^3)$.
\begin{enumerate}
  \item Que cela implique-t-il avec 500\,000 observations ?
  \item Citer deux alternatives pratiques.
  \item Le noyau RBF change-t-il cet ordre de grandeur ?
\end{enumerate}

\newpage

\section*{TD 7 --- Apprentissage non supervisé}
\addcontentsline{toc}{section}{TD 7 --- Apprentissage non supervisé}

\subsection*{Exercice 1 --- Une itération de $k$-means à la main \hfill \normalfont\footnotesize (12 min)}
Quatre points : $P_1=(1,1)$, $P_2=(2,1)$, $P_3=(4,5)$, $P_4=(5,4)$.
Centres initiaux : $C_1=(1,1)$, $C_2=(5,4)$.
\begin{enumerate}
  \item Calculer les distances au carré de chaque point aux deux centres.
  \item Affecter chaque point au centre le plus proche.
  \item Recalculer les deux centres.
  \item Les affectations changeraient-elles à l'itération suivante ?
\end{enumerate}

\subsection*{Exercice 2 --- Choisir $k$ \hfill \normalfont\footnotesize (8 min)}
On trace l'inertie intra-classe en fonction de $k$ : elle décroît de 120 à 40
entre $k=1$ et $k=3$, puis de 40 à 34 entre $k=3$ et $k=8$.
\begin{enumerate}
  \item Où se situe le « coude » ?
  \item Pourquoi l'inertie ne permet-elle pas de choisir $k$ seule ?
  \item Quel autre critère utiliser, et que mesure-t-il ?
\end{enumerate}

\subsection*{Exercice 3 --- ACP \hfill \normalfont\footnotesize (8 min)}
Une ACP sur 4 variables donne des variances expliquées de 60 \%, 25 \%, 10 \%
et 5 \%.
\begin{enumerate}
  \item Combien de composantes faut-il pour dépasser 85 \% ?
  \item Que contient la quatrième composante ?
  \item Pourquoi l'ACP impose-t-elle de standardiser les variables ?
\end{enumerate}

\subsection*{Exercice 4 --- DBSCAN ou $k$-means ? \hfill \normalfont\footnotesize (8 min)}
Indiquer l'algorithme adapté :
\begin{enumerate}
  \item trois groupes de tailles comparables, bien séparés ;
  \item des groupes en forme de croissant ;
  \item des points isolés que l'on veut signaler comme bruit.
\end{enumerate}

\subsection*{Exercice 5 --- Détection d'anomalies \hfill \normalfont\footnotesize (6 min)}
\begin{enumerate}
  \item Pourquoi ne peut-on pas entraîner un classifieur supervisé quand les
        anomalies sont très rares ?
  \item Que fait l'Isolation Forest, en une phrase ?
  \item Quelle métrique utiliser pour évaluer un détecteur d'anomalies ?
\end{enumerate}


\subsection*{Exercice 6 --- Calculer une silhouette \hfill \normalfont\footnotesize (8 min)}
Pour un point : distance moyenne aux points de son groupe $a=2$ ; distance
moyenne au groupe voisin $b=6$.
\begin{enumerate}
  \item Calculer la silhouette $(b-a)/\max(a,b)$.
  \item Que vaut-elle si $a=5$ et $b=5$ ? Comment l'interpréter ?
  \item Quelle est la plage de valeurs possibles ?
\end{enumerate}

\subsection*{Exercice 7 --- Reconstruire après une ACP \hfill \normalfont\footnotesize (7 min)}
On garde une seule composante expliquant 60 \% de la variance.
\begin{enumerate}
  \item Que perd-on en reconstruisant les données d'origine ?
  \item Pourquoi 60 \% peut-il suffire pour visualiser, mais pas pour décider ?
  \item Que mesure la variance expliquée cumulée ?
\end{enumerate}

\subsection*{Exercice 8 --- Minima locaux de $k$-means \hfill \normalfont\footnotesize (7 min)}
Deux initialisations différentes donnent deux partitions différentes.
\begin{enumerate}
  \item Pourquoi l'algorithme dépend-il de l'initialisation ?
  \item Que fait l'initialisation $k$-means++ ?
  \item Comment rendre le résultat plus stable ?
\end{enumerate}

\subsection*{Exercice 9 --- Choisir un seuil d'anomalie \hfill \normalfont\footnotesize (7 min)}
L'Isolation Forest fournit un score continu.
\begin{enumerate}
  \item Pourquoi ne pas prendre un seuil arbitraire à 0,5 ?
  \item Quel coût comparer pour choisir le seuil ?
  \item Quelle métrique évaluer en l'absence d'étiquettes ?
\end{enumerate}

\newpage

\section*{TD 8 --- Fiabilité, éthique et mise en production}
\addcontentsline{toc}{section}{TD 8 --- Fiabilité, éthique et mise en production}

\subsection*{Exercice 1 --- Auditer un rapport \hfill \normalfont\footnotesize (12 min)}
Un rapport annonce : « 98 \% d'exactitude pour prédire l'abandon scolaire ».
On y lit que le modèle utilise « le nombre de semestres validés » et « la
décision d'exclusion », et que le découpage est aléatoire sur 5 promotions.
\begin{enumerate}
  \item Relever la fuite métier.
  \item Relever le problème de protocole.
  \item Quelles deux exigences poser avant de valider ce travail ?
\end{enumerate}

\subsection*{Exercice 2 --- Biais et équité \hfill \normalfont\footnotesize (8 min)}
Un modèle de sélection de candidats est entraîné sur 10 ans d'historique, où
les femmes étaient sous-représentées dans les postes techniques.
\begin{enumerate}
  \item Pourquoi le modèle risque-t-il de reproduire ce déséquilibre ?
  \item Suffit-il de retirer la variable « sexe » ? Justifier.
  \item Proposer deux mesures de contrôle.
\end{enumerate}

\subsection*{Exercice 3 --- Dérive des données \hfill \normalfont\footnotesize (8 min)}
Un modèle de prévision de la demande, entraîné avant une crise économique,
perd 15 points de performance six mois plus tard.
\begin{enumerate}
  \item Nommer le phénomène.
  \item Comment le détecter avant que les utilisateurs ne se plaignent ?
  \item Que faire : réentraîner, recalibrer ou changer de modèle ?
\end{enumerate}

\subsection*{Exercice 4 --- Interprétable ou performant ? \hfill \normalfont\footnotesize (8 min)}
Deux modèles : une régression logistique (AUC 0,81, coefficients lisibles) et
un gradient boosting (AUC 0,86, opaque).
\begin{enumerate}
  \item Dans quel contexte choisir le premier ?
  \item Dans quel contexte le second ?
  \item Comment réduire l'écart entre les deux ?
\end{enumerate}

\subsection*{Exercice 5 --- Trois limites honnêtes \hfill \normalfont\footnotesize (6 min)}
Rédiger trois limites d'un modèle de prédiction de la réussite étudiante,
chacune en une phrase, en distinguant : les données, la méthode, l'usage.


\subsection*{Exercice 6 --- Fuite temporelle \hfill \normalfont\footnotesize (7 min)}
Un modèle de prévision de la demande est entraîné sur 2019--2023 et testé sur
des mois tirés au hasard dans cette période.
\begin{enumerate}
  \item Expliquer pourquoi ce protocole est invalide.
  \item Proposer un découpage correct.
  \item Quel effet sur le score annoncé ?
\end{enumerate}

\subsection*{Exercice 7 --- Variables proxy \hfill \normalfont\footnotesize (7 min)}
On retire la variable « sexe » d'un modèle de recrutement.
\begin{enumerate}
  \item Pourquoi le biais peut-il subsister ?
  \item Citer deux variables pouvant servir de proxy.
  \item Comment vérifier l'équité du modèle après entraînement ?
\end{enumerate}

\subsection*{Exercice 8 --- Surveillance en production \hfill \normalfont\footnotesize (7 min)}
\begin{enumerate}
  \item Citer quatre indicateurs à surveiller après déploiement.
  \item Lesquels se détectent sans attendre les étiquettes réelles ?
  \item Que faire si la distribution d'une variable d'entrée se décale ?
\end{enumerate}

\subsection*{Exercice 9 --- Fiche modèle \hfill \normalfont\footnotesize (6 min)}
Une \emph{model card} documente un modèle. Citer cinq rubriques qu'elle doit
contenir et, pour chacune, une phrase sur son utilité.

\newpage

\section*{TD 9 --- Synthèse : sujet de projet}
\addcontentsline{toc}{section}{TD 9 --- Synthèse : sujet de projet}

\subsection*{Exercice unique --- Cadrer un projet de bout en bout \hfill \normalfont\footnotesize (45 min)}
Une école veut prédire quels étudiants risquent d'échouer, pour proposer un
tutorat dès le premier mois. On dispose de cinq promotions, de la moyenne au
premier semestre, du nombre d'absences, de la filière et du résultat final.

\begin{enumerate}
  \item \textbf{Données.} Quelles variables sont disponibles au moment de la
        décision ? Lesquelles sont postérieures et doivent être écartées ?
  \item \textbf{Découpage.} Quel schéma de validation ? Justifier en une phrase.
  \item \textbf{Baseline.} Quelle baseline proposez-vous, et quel score
        obtiendrait-elle approximativement si 70 \% des étudiants réussissent ?
  \item \textbf{Métrique.} L'école veut éviter de rater un étudiant en
        difficulté. Quelle métrique privilégier, et quelles en sont les
        conséquences ?
  \item \textbf{Fuite.} Citer deux fuites possibles dans ce projet précis.
  \item \textbf{Usage.} Rédiger une phrase décrivant ce que l'école fera du
        score, et une phrase décrivant ce qu'elle ne doit pas en faire.
\end{enumerate}

\vspace{3mm}
{\footnotesize\color{gris}
Ce TD reprend la structure du projet intégrateur : cadrage, découpage,
baseline, métrique, fuite, usage. Un groupe qui répond correctement à ces six
questions a fait l'essentiel du travail de conception.}


\subsection*{Exercice 2 --- Coût et faisabilité \hfill \normalfont\footnotesize (10 min)}
\begin{enumerate}
  \item Estimer le volume de travail : collecte, annotation, modélisation,
        application.
  \item Quel est le risque principal du projet, et comment le réduire tôt ?
  \item Quel critère d'arrêt (go / no-go) poser au jalon 1 ?
\end{enumerate}

\subsection*{Exercice 3 --- Résultats attendus \hfill \normalfont\footnotesize (10 min)}
\begin{enumerate}
  \item Quel gain l'école peut-elle espérer, en points de rappel ?
  \item Quelle baseline faut-il battre pour que le projet ait un intérêt ?
  \item Comment mesurer que l'outil est réellement utilisé ?
\end{enumerate}

\subsection*{Exercice 4 --- Soutenance \hfill \normalfont\footnotesize (8 min)}
Préparer trois réponses courtes :
\begin{enumerate}
  \item « Pourquoi votre score est-il crédible ? »
  \item « Sur quels étudiants vous trompez-vous le plus ? »
  \item « Que fera l'école de vos prédictions ? »
\end{enumerate}

\newpage


\end{document}
